\documentclass[12pt,a4paper]{article}

\usepackage{palatino}
\usepackage{amsmath,amssymb,amsthm}
\usepackage{geometry}
\usepackage{microtype}
\usepackage{hyperref}
\usepackage{setspace}
\usepackage{csquotes}
\usepackage{titlesec}
\usepackage{abstract}
\usepackage{mathtools}

\geometry{
  top=1.25in,
  bottom=1.25in,
  left=1.25in,
  right=1.25in
}

\onehalfspacing

\hypersetup{
  colorlinks=true,
  linkcolor=black,
  citecolor=black,
  urlcolor=black
}

\titleformat{\section}{\large\bfseries}{\thesection.}{0.75em}{}
\titleformat{\subsection}{\normalsize\bfseries}{\thesubsection.}{0.75em}{}

\theoremstyle{plain}
\newtheorem{theorem}{Theorem}
\newtheorem{proposition}{Proposition}
\newtheorem{lemma}{Lemma}
\newtheorem{corollary}{Corollary}

\theoremstyle{definition}
\newtheorem{definition}{Definition}
\newtheorem{example}{Example}

\theoremstyle{remark}
\newtheorem{remark}{Remark}

\title{\textbf{Reality Is What Can Be Reached:\\
Reachability, Identification, and the Limits of Abstraction}}

\author{Flyxion}

\date{June 2026}

\begin{document}

\maketitle

\begin{abstract}
\noindent
Modern science, statistics, artificial intelligence, and philosophy largely
inherit a common ontological assumption: reality consists fundamentally of
objects, states, facts, and entities whose properties are subsequently
described, predicted, compressed, classified, and optimized. This essay
argues for a reversal of that picture. The primitive structure of reality is
not objecthood but reachability: the pattern of futures accessible from a
given state. Objects, representations, abstractions, scientific laws,
optimization procedures, and predictive models emerge as secondary
constructions imposed upon this deeper accessibility geometry.

From this perspective, every ontology becomes a theory of identification,
every representation becomes a projection, every abstraction becomes a form of
controlled collapse, and every act of compression becomes an act of selective
forgetting. The central question is therefore not whether distinctions are
removed, but which distinctions may be removed without altering the future
structure of the system under consideration. This criterion is formalized
through the concept of admissibility, understood as the preservation of
reachable futures under projection.

The essay develops a unified framework connecting abstraction, compression,
prediction, optimization, scientific modeling, and artificial intelligence
through the geometry of accessible futures. It argues that prediction,
compression, statistical significance, and optimization are not foundational
principles but downstream instruments whose validity depends upon a prior
condition: the preservation of reachability structure. A system may be highly
predictive, highly compressive, statistically successful, and optimally
controlled while simultaneously destroying the distinctions that determine
which futures remain possible.

The resulting inversion has significant consequences for contemporary debates
in machine learning, control theory, epistemology, and the philosophy of
science. Intelligence is reinterpreted not as the discovery of regularities,
the maximization of prediction accuracy, or the compression of experience,
but as the preservation of future possibility. The deepest question becomes
not what is true, predictable, or compressible, but what cannot be forgotten
without destroying the future.
\end{abstract}

\bigskip

%%% SECTION 1
\section{The Crack in the Bridge}

Consider a bridge carrying thousands of vehicles each day. For years it
behaves exactly as expected. Its load-bearing characteristics remain within
acceptable engineering tolerances. Traffic flows normally. Weather patterns
remain unexceptional. Every available metric suggests stability. The bridge,
from the perspective of any standard observational framework, is functioning
as designed.

Somewhere within the structure, however, a microscopic crack begins to form.

Initially the crack contributes almost nothing to any conventional measure of
importance. It explains essentially none of the bridge's observed behavior.
It has negligible influence on average load calculations. It contributes
almost nothing to predictive error. Under many forms of statistical modeling
it would be treated as noise and, in sufficiently compressed representations,
would likely vanish entirely---swallowed by the equivalence classes of a
model trained to minimize description length or expected forecast loss.

Yet the crack matters, and it matters enormously.

The reason it matters has little to do with its current magnitude and
everything to do with its relationship to the future. Although physically
small, it occupies a position within the structure from which possible
futures diverge. One future contains a functioning bridge and normal daily
use. Another contains progressive structural degradation and costly repair.
Another contains catastrophic failure. The significance of the crack does not
arise from its contribution to present observations but from its capacity to
alter what futures remain accessible from the current state of the system.

This distinction reveals a tension that appears repeatedly throughout
science, engineering, and artificial intelligence. Many of the most
influential evaluation frameworks assess importance according to present
contribution. Statistical methods measure variance explained. Predictive
systems measure reductions in expected error. Compression methods measure
reductions in description length. Optimization procedures evaluate changes in
objective function values. In each case, the significance of a feature is
assessed through its relationship to some quantity defined relative to
present observations or present performance.

The crack exposes a limitation shared by all such approaches. What
determines the significance of the crack is not the amount of information it
contributes to current observations, nor the degree to which its inclusion
improves prediction, nor the resistance it offers to compression. What
determines its significance is its influence on future accessibility. A
detail may contribute almost nothing to present description while
simultaneously controlling the boundary between radically different futures.

The same phenomenon appears across many domains. A latent software fault may
remain dormant for years before triggering catastrophic system failure at
scale. A single mutation may redirect the evolutionary trajectory of an
entire species across geological time. A pathogen establishing itself in a
small isolated population may transform into the beginning of an epidemic. A
small but legally ambiguous financial liability may remain invisible within
ordinary accounting metrics until economic conditions change and suddenly
expose its consequences. A subtle misunderstanding between political actors
may persist unnoticed for years until it becomes the catalyst for large-scale
conflict.

In every case the critical feature is the same. The present statistical
footprint of the state is small. Its future consequences are enormous. The
problem is not merely practical. It is, at its root, ontological.

Modern intellectual culture tends to assume that reality consists primarily
of objects, states, facts, and entities. Once these objects have been
identified, the task of science becomes describing them, predicting their
behavior, compressing their structure, and optimizing their interactions. The
crack appears, within this picture, as merely another property of another
object---a physical quantity to be measured, predicted, and managed along
with all the others.

Yet the example suggests a different possibility. Perhaps the significance of
a state does not derive primarily from what it is but from what it permits.
Perhaps the fundamental structure of reality is not objecthood but
accessibility. Perhaps what matters most about any state is not its present
description but the pattern of futures that remain reachable from it.

If this inversion is correct, then many familiar intellectual priorities must
be reconsidered. Prediction would no longer be fundamental but derivative.
Compression would no longer define understanding. Optimization would become
a special case of navigation within a pre-existing structure of possibility.
Representation would become a problem of preserving future accessibility
rather than merely preserving information.

The crack therefore serves as more than an engineering example. It reveals a
general principle, which may be stated as follows: the features that matter
most about a system are not necessarily those that contribute most strongly
to present description, prediction, compression, or control. They are the
features whose removal would alter the structure of reachable futures. The
remainder of this essay develops the consequences of taking that principle
seriously as a foundational commitment rather than a useful heuristic.

%%% SECTION 2
\section{Reality Is What Can Be Reached}

The example of the crack motivates a shift in ontological priority that runs
against several deeply entrenched intellectual habits. Ordinarily, reality is
assumed to consist of objects possessing properties. A bridge contains steel,
concrete, bolts, supports, and defects. A biological organism contains cells,
organs, and molecules. A physical system contains particles, fields, and
interactions. The task of scientific explanation is then understood as the
discovery of laws governing the behavior of these objects and their
properties. This picture is so familiar that it often appears self-evident,
as though the priority of objects were itself a feature of the world rather
than a theoretical commitment about what to take as primitive.

Yet this picture contains an important assumption that is rarely made
explicit. It assumes that objects are ontologically primary and that
possibilities emerge from the interactions among those objects. The world is
first populated with things; only afterwards do futures appear. The crack, on
this view, is a defect in a structural member, and its significance is a
downstream consequence of what that defect eventually does to the properties
of the surrounding material.

The crack example points consistently in the opposite direction. What made
the crack important was not its present physical description. It was not
especially large, it did not dominate any observable metric, and it was not
responsible for most of the bridge's behavior. The significance of the crack
derived from the way it altered future accessibility. Its importance arose
from its location within a network of possible transformations---from the
fact that it stood at a branching point in the space of reachable futures.

This observation motivates a different starting point. Rather than beginning
with objects and deriving futures, one may begin with futures and derive
objects. Instead of treating reality as a collection of entities that happen
to possess possibilities, one may treat reality as a structure of
possibilities from which entities emerge as derivative constructs. Objects,
on this view, are not the bedrock of reality but are stabilized patterns
within a deeper geometry of accessible transformations---regions of the
reachability structure that persist coherently across a range of conditions
and timescales.

The inversion can be expressed schematically. The traditional picture assumes
a hierarchy of the form
\[
  \text{Reality}
  \;\rightarrow\;
  \text{Objects}
  \;\rightarrow\;
  \text{Relations}
  \;\rightarrow\;
  \text{Futures},
\]
where objects come first and their relational dynamics subsequently give rise
to the space of possible futures. The reachability picture proposed here
inverts this ordering:
\[
  \text{Reality}
  \;\rightarrow\;
  \text{Reachability Structure}
  \;\rightarrow\;
  \text{Objects}.
\]
Objects become secondary constructions. They emerge as relatively stable
regions within a deeper geometry of accessible transformations, and their
significance derives from the role they play in that geometry rather than
the other way around.

This shift may initially appear radical, but versions of it already appear
throughout several domains of science and mathematics. In dynamical systems
theory, the long-term behavior of a system is routinely analyzed through
attractors, basins of attraction, bifurcation diagrams, and invariant
manifolds rather than through isolated states. The state of the system
matters less than the qualitative structure of the space of trajectories
emanating from it. In ecology, the significance of a species is frequently
determined by the role it plays in sustaining accessible pathways through an
ecosystem: a keystone species matters not primarily because of what it is
but because of what its removal forecloses. In economics, the value of an
asset derives not merely from present ownership but from the opportunities
that ownership makes available---an option, for instance, derives its value
entirely from the accessibility of future states it preserves. In computer
science, abstract data types are often defined operationally through the
transformations they permit rather than through their internal constitution
alone, making accessibility part of what the structure fundamentally is.

To formalize the reachability framework, let $x$ denote a state of some
system and let $\mathcal{R}(x)$ denote the set of futures reachable from
that state under an appropriate class of transformations or policies. The
object of primary theoretical interest is no longer the state itself but the
structure of $\mathcal{R}(x)$ as a function across the state space. The
meaning of a state, in this sense, becomes inseparable from its future
horizon: two states that appear superficially similar may differ profoundly
if the futures accessible from them differ, while states that appear
physically distinct may be equivalent for many purposes if they permit the
same future transformations. What matters is not merely what the state is
but what the state allows.

This perspective alters the interpretation of knowledge itself. Knowledge is
commonly understood as accurate description: one possesses knowledge to the
extent that one's internal representations correspond faithfully to the
external world. Yet from a reachability perspective, the practical value of a
representation depends less on descriptive fidelity than on its ability to
track future accessibility. A representation succeeds not when it resembles
the world but when it preserves the distinctions that determine which futures
remain possible. Two representations may differ dramatically in their surface
properties while remaining equivalent from the perspective of reachability
preservation, and two representations may appear similar while differing
profoundly in the futures they allow one to navigate.

The same inversion affects prediction. Prediction is often treated as the
defining feature of intelligence and the primary criterion for evaluating
scientific theories and artificial systems alike. Machine learning systems
are trained to minimize predictive error. World models are evaluated
according to their ability to forecast future states. Scientific theories are
judged by predictive success. Yet prediction already presupposes accessibility
structure. To predict is to make claims about possible futures; to plan is to
select among possible futures; to control is to guide movement toward possible
futures; and to optimize is to search among possible futures. Even safety is
defined through the distinction between acceptable and unacceptable futures.
In every case, the concept of a future is logically prior to the operation
being performed upon it.

This observation reveals a subtle asymmetry. Prediction, planning,
optimization, and control are often treated as foundational concepts, yet
they all depend upon a more primitive notion: the structure of what can be
reached. Without some prior notion of accessibility, there is nothing to
predict, nothing to optimize, and nothing to control. Reachability is
therefore not introduced here as an alternative objective alongside prediction
or optimization, nor as a constraint to be satisfied alongside other
requirements. It is introduced as a prior condition for those activities to
make sense at all.

A clarification is warranted here to forestall a possible misreading. The
reachability thesis does not assert that objects are unreal or that physical
properties are irrelevant. Steel possesses specific tensile properties; cracks
propagate according to known mechanical laws; bridges fail when load exceeds
capacity. Object-level descriptions explain \emph{why} futures diverge.
Reachability explains \emph{why that divergence matters}. The two levels of
description are complementary and mutually dependent, not competing. The
proposal concerns explanatory priority for a specific purpose: determining
which distinctions deserve preservation under abstraction. For that purpose,
the accessibility geometry provides the more general and directly applicable
criterion, because it is the geometry within which planning, control, and
safety must ultimately operate.

Furthermore, reachability is never absolute or free-floating. It is always
defined relative to a specified dynamical system $D$ encoding the laws,
constraints, and permitted transformations of the domain under consideration.
Writing this dependence explicitly, the reachable future set from state $x$
under dynamics $D$ is
\[
  \mathcal{R}_D(x)
  =
  \{y \in \mathcal{X} : x \rightsquigarrow_D y\},
\]
where $x \rightsquigarrow_D y$ denotes reachability under the transition
structure of $D$. The object-level description---material properties,
governing equations, boundary conditions---is precisely what determines $D$.
The present argument is not that reachability precedes dynamics in any
metaphysically loaded sense; it is that once the dynamical structure has been
specified, the accessibility geometry $\mathcal{R}_D$ becomes the natural
criterion for evaluating which representational collapses are permissible.
Object properties determine the geometry. Admissibility evaluates projections
relative to that geometry.

%%% SECTION 3
\section{Why Reachability Comes First}

The claim that reachability is ontologically prior to prediction, planning,
optimization, and control may initially appear counterintuitive. Contemporary
scientific and engineering practice often treats these latter concepts as
foundational. Entire disciplines are organized around the prediction of future
states, the optimization of objective functions, and the control of dynamical
systems. It is therefore natural to ask why reachability should be granted a
more primitive status rather than being regarded as simply one concept among
others.

The answer lies not in the practical importance of these activities but in
their logical dependence upon a prior structure that they presuppose but
cannot themselves provide.

Consider prediction. To predict is to make a statement about a future state
of some system. Even the simplest prediction presupposes that the future
state being discussed belongs to a space of possible futures that already
possesses determinate structure. A prediction may be correct or incorrect,
but before it can be either, there must exist some underlying geometry of
possibilities that determines which futures are available and which are not.
The act of prediction therefore presupposes an accessibility structure that
the prediction itself does not create and cannot modify.

Planning exhibits the same dependency in a particularly transparent form. A
plan is a sequence of actions intended to move a system from one state to
another. The very concept of planning requires that multiple futures be
available for consideration: a planner evaluates alternative trajectories,
compares their consequences, and selects among them. Yet none of these
operations determine which trajectories exist in the first place. Planning
operates within a space of possibilities whose structure is already given,
and the quality of the plan depends entirely upon the admissibility of the
representation through which that space is perceived.

Control theory provides an especially clear example. A controller seeks to
guide a system toward some desired region of state space. Classical control
theory is extraordinarily successful precisely because it typically assumes
that the relevant state variables have already been identified and that the
transition structure of the system is sufficiently understood. Once those
assumptions are granted, the controller can exploit the geometry of the state
space to produce desired behavior. What the controller does not do is
generate that geometry. The geometry exists prior to the controller. The
controller navigates it.

Optimization reveals the same asymmetry. Optimization is frequently treated
as a foundational explanatory principle: contemporary machine learning,
economics, and decision theory often describe intelligent behavior as the
maximization of some objective function over a space of possible states. Yet
optimization is incapable of defining that space. An optimizer can search
among alternatives, rank futures by objective value, and identify maxima or
minima. It cannot determine what counts as an alternative in the first place.
An optimizer requires a geometry of possibilities to exist before it can
begin operating, yet it takes no responsibility for the admissibility of that
geometry.

This distinction is easy to overlook because optimization procedures often
appear enormously powerful. Given a sufficiently rich state space and a
sufficiently expressive objective function, optimization can produce
remarkably complex and adaptive behavior. The success of the optimizer,
however, depends entirely upon the structure of the space within which it
operates. An optimizer searching the wrong geometry may be perfectly rational
in a formal sense while failing catastrophically in practice.

The familiar observation that a system is ``optimizing the wrong objective''
captures only part of the problem. A deeper possibility exists and is in many
ways more troubling: a system may be optimizing the correct objective within
the wrong representation. In such a case the failure does not arise from the
objective function itself but from the geometry of distinctions available to
the optimizer. The objective is fine. The map is wrong. And crucially, the
optimizer has no means of detecting this from within its own representational
world.

Safety illustrates the dependency upon reachability with particular clarity.
Safety is often discussed as though it were an additional objective or
constraint that can be imposed upon an otherwise capable system. One
introduces guardrails, penalties, alignment criteria, or constraint
satisfaction requirements intended to prevent undesirable behavior. Such
measures are undeniably important, but they presuppose that the system can
distinguish safe states from unsafe states in the first place. If the
underlying representation has already collapsed those states into a single
equivalence class, no downstream safety mechanism can recover the distinction.
The distinction has been destroyed at the level of projection, and guardrails
operating above that level inherit the blindness without being able to
compensate for it.

The logical structure of these observations can be summarized as a chain of
dependencies. Prediction presupposes a space of accessible futures.
Planning presupposes that the space of futures has been correctly
represented. Optimization presupposes a geometry within which the search
occurs. Control presupposes a representation in which relevant distinctions
survive. Safety presupposes that catastrophic and acceptable states remain
distinguishable. Each activity therefore depends upon a prior condition that
none of them supply: the preservation of reachability structure under
abstraction.

Reachability thus occupies a different logical category from the activities
that depend upon it. It is not another objective alongside prediction or
optimization, to be balanced against other requirements in some
multi-objective framework. It is the condition under which prediction,
optimization, planning, control, and safety become coherent as activities at
all. Understanding why this is so requires examining what happens when
systems form representations---how they identify states as equivalent and
thereby decide, always implicitly and often irreversibly, which distinctions
survive.

%%% SECTION 4
\section{Every Ontology Is a Theory of Identification}

If reachability is primary, then the central problem of knowledge changes
substantially. The question is no longer simply how to describe the world but
how to partition it. Every act of understanding requires determining which
differences matter and which may be safely ignored. Every representation,
category, concept, scientific law, and model performs this operation, and
performs it necessarily---for no finite cognitive or computational system can
preserve every distinction present in the world it inhabits.

This observation appears trivial at first glance. Human beings constantly
identify things as equivalent. Individual trees are identified as instances
of the same species. Physically distinct chairs are subsumed under the same
category. Successive versions of a nation, a language, or a person are
treated as manifestations of a single continuing entity despite enormous
internal change. Scientific theories likewise identify apparently different
phenomena as instances of common principles. The history of science is filled
with such unifications: heat and molecular motion, electricity and magnetism,
space and time, mass and energy. Each represents an identification that
collapsed a previously maintained distinction.

Yet beneath these familiar operations lies a deeper question. What justifies
any particular identification? The traditional answer appeals to similarity:
two entities are treated as equivalent because they resemble one another in
relevant respects. This answer, however, merely relocates the problem.
Similarity itself requires a criterion. To determine whether two states are
similar one must already know which features matter and which do not. The act
of identification cannot be grounded entirely in resemblance because
resemblance presupposes that relevant dimensions of comparison have already
been selected, and that selection is itself an act of identification. The
problem is therefore circular if grounded only in similarity.

The reachability framework provides an alternative foundation. To identify
two states as equivalent is to collapse a distinction between them. Every
category, concept, representation, and theory performs such a collapse.
Every ontology therefore enacts a theory of permissible identification: it
specifies which distinctions may be ignored without loss of the explanatory
or operational power that matters for the domain in question. Ontology, on
this view, is not primarily a catalogue of what exists. It is a theory of
what can be treated as the same.

Consider the concept of a species. Individual organisms differ in countless
ways---different genetic sequences, developmental histories, behaviors,
morphological features, and environmental interactions. The category of
species identifies many of these organisms as equivalent despite those
differences. The category is useful precisely because it ignores distinctions
that are irrelevant for many biological purposes while preserving the
distinctions that matter for reproduction, inheritance, and evolutionary
dynamics. A similar analysis applies to the concept of a physical object.
A chair remains a chair despite scratches, repainting, incremental repairs,
and gradual material replacement over time. The concept functions by
collapsing numerous physical differences into a single persisting identity.
Without such collapse, ordinary perception and action would become impossible.

Scientific laws operate analogously. Newtonian mechanics identifies countless
distinct motions as instances of common dynamical principles, collapsing the
enormous diversity of trajectories in the physical world into a small
collection of equations. Thermodynamics ignores microscopic molecular details
to describe macroscopic behavior, treating the uncountable individual
trajectories of gas molecules as equivalent so long as they share the same
macroscopic state. Fluid mechanics treats vast collections of individual
particles as continuous media, collapsing distinctions among molecular
configurations that leave bulk flow behavior unchanged. These abstractions
succeed not because they preserve all information but because they preserve
the right information for their domain of application.

The success of science therefore depends not on avoiding identification but
on performing the correct identifications. A scientific theory that
maintained every microscopic distinction would be computationally intractable
and explanatorily useless. What matters is the criterion by which
identifications are made.

Representations in artificial intelligence perform the same operation. A
learned representation maps many distinct physical states into a smaller
latent space, declaring states equivalent whose latent codes coincide. The
representation succeeds precisely to the extent that it identifies states
that may safely be treated as equivalent for the purposes of prediction,
planning, or control. The language of representation learning often describes
this process in terms of feature extraction, dimensionality reduction,
abstraction, or compression, but ontologically the operation is the same: the
representation decides which distinctions survive and which do not.

This perspective reveals a hidden continuity between seemingly unrelated
intellectual activities. The philosopher constructing an ontology, the
scientist formulating a law, the engineer building a model, and the machine
learning system learning a latent representation are all solving variations
of the same problem. Each must determine which distinctions are worth
preserving and which may be collapsed without harmful consequence.

The challenge is that collapse is irreversible in practice. Once two states
have been identified as equivalent within a representation, every subsequent
operation inherits that decision. Predictions, optimizations, classifications,
and plans all operate on the resulting equivalence classes. They cannot
easily recover distinctions that have already been removed by the
representation upon which they depend. The decision about which distinctions
to preserve is therefore among the most consequential decisions a cognitive or
computational system makes, and it is often made implicitly, in the structure
of representation choices, rather than explicitly and deliberately.

The significance of this observation becomes clearer when viewed through the
lens of reachability. If reality is fundamentally a geometry of accessible
futures, then the adequacy of an identification cannot be judged solely by
similarity, predictive success, or compression efficiency. A more demanding
criterion is required. Two states may be safely identified only if the
collapse preserves the future structure of the system. If the identification
destroys distinctions that alter what futures remain accessible, then the
collapse is not merely imperfect---it is inadmissible.

%%% SECTION 5
\section{The Projection Problem}

The preceding discussion established that every ontology performs acts of
identification. Categories, scientific laws, abstractions, and
representations all reduce complexity by treating certain distinctions as
irrelevant. The critical question is therefore not whether collapse occurs
but how it is performed, and whether the collapses that occur are warranted.
To answer this question with greater precision, it is useful to formalize the
relationship between a system and its representation.

Let $\mathcal{X}$ denote a state space and let $\mathcal{M}$ denote a
representation space. A representation can be understood as a mapping
\[
  \pi : \mathcal{X} \rightarrow \mathcal{M}.
\]
The notation is intentionally general. The representation $\pi$ may
correspond to a scientific model, a statistical summary, a learned latent
embedding, a conceptual category, a digital twin, a control state, or any
other structure intended to stand in for the original system.

The essential feature of such mappings is that they are almost never
injective. Distinct states in $\mathcal{X}$ are mapped to the same
representation in $\mathcal{M}$. Formally, for any non-trivial projection
there exist states $x_1, x_2 \in \mathcal{X}$ such that
\[
  x_1 \neq x_2
  \qquad\text{yet}\qquad
  \pi(x_1) = \pi(x_2).
\]
Whenever this occurs, an identification has been performed. The representation
has declared the distinction between $x_1$ and $x_2$ irrelevant---collapsed
it into a shared representational point from which neither element can be
independently recovered.

The term \emph{projection} is more precise than \emph{abstraction} or
\emph{representation} because it emphasizes what these mappings actually do.
A projection does not merely preserve some information while discarding
other information. It actively eliminates distinctions by imposing an
equivalence relation on the original state space. The equivalence classes
generated by $\pi$---the sets $[x]_\pi = \{x' : \pi(x') = \pi(x)\}$---
define what the representation can and cannot express. Every claim made within
$\mathcal{M}$ is implicitly a claim about equivalence classes rather than
individual states, and the consequences of actions taken on the basis of that
representation are therefore consequences about sets of states rather than
the specific state actually occupied.

This observation immediately reveals an asymmetry that is underappreciated in
most discussions of representation learning. Much of the contemporary
discussion focuses on what representations capture. Researchers ask whether a
latent space contains information about depth, object identity, causal
structure, planning relevance, or commonsense relations. Such questions are
important, but they overlook an equally significant question that is in many
ways logically prior.

Every representation also defines what has been forgotten.

The omissions are not incidental artifacts of imperfect learning. They are
the price of abstraction itself, and they are paid by necessity with every
representational choice. The fact that they are often invisible---that
successful abstractions appear natural in retrospect---does not diminish their
importance. It merely makes them harder to examine critically.

Traditional evaluation frameworks typically address the projection problem
indirectly. Statistical models evaluate projections according to predictive
performance on held-out data. Compression methods evaluate them according to
description length. Information-theoretic approaches evaluate them according
to preserved variance or mutual information. Optimization frameworks evaluate
them according to downstream task success in controlled environments.

All of these criteria have practical value, and none of them should be
abandoned. However, none directly addresses the question that admissibility
poses. A projection may perform extremely well according to any of these
measures while still destroying distinctions that alter the future structure
of the system. The reason is straightforward: prediction, compression, and
optimization evaluate properties that are internal to the representation, or
relative to some training distribution. They do not directly evaluate what
futures have been lost through the act of projection itself, because the
lost futures are precisely what is no longer visible from within the
representational space.

The crack in the bridge again provides a useful illustration. Suppose a
representation collapses two bridge states into a single latent description
because their observable behavior is nearly identical across the training
distribution. The resulting projection may improve compression, reduce
prediction error, and simplify the downstream control problem. Yet if one
state contains a microscopic structural defect and the other does not, the
projection has destroyed a distinction that separates radically different
future trajectories. The representation has become simpler; the future has
become less visible. And crucially, the standard evaluation metrics will
record the projection as having succeeded, because the failure is located
precisely in the region of the state space that has been collapsed out of
representational existence.

The central intuition motivating this essay can now be stated more precisely.
Representational adequacy cannot be assessed solely through predictive or
informational criteria because those criteria evaluate performance within a
given geometry of distinctions, while the most important failures occur when
the geometry itself has been incorrectly specified. A projection is successful
not merely when it preserves information in the aggregate but when it
preserves the distinctions necessary to maintain the future structure of the
system---when, in other words, it is admissible.

%%% SECTION 6
\section{Admissibility}

The projection problem reveals that a representation may preserve large
quantities of information, support accurate prediction, enable efficient
compression, and facilitate successful optimization while still destroying
distinctions that fundamentally alter the future structure of the system. If
reachability is primary, then a different criterion is needed to distinguish
productive abstraction from destructive collapse. This criterion is
admissibility.

The concept begins with a structural observation about every state. Let
$x \in \mathcal{X}$ denote a state of some system. Associated with $x$ is
a collection of futures reachable through the allowable transformations of
that system. Denote this collection by $\mathcal{R}(x)$. The precise
meaning of reachability depends upon context: in a control problem,
$\mathcal{R}(x)$ may represent the states accessible through some class of
control policies; in a biological system, it may represent possible future
evolutionary trajectories; in an economic system, it may represent attainable
economic configurations under feasible sequences of decisions. The formal
details differ across domains, but the underlying idea remains consistent:
every state determines a horizon of futures, and that horizon is among the
most fundamental properties of the state from the perspective of any system
that must act within the world.

\begin{definition}[Admissible Projection]
Let $\pi : \mathcal{X} \rightarrow \mathcal{M}$ be a projection from a
state space to a representation space, and let $\mathcal{R}(x)$ denote the
reachable future set from state $x$. The projection $\pi$ is said to be
\emph{admissible} if for all $x_1, x_2 \in \mathcal{X}$,
\[
  \pi(x_1) = \pi(x_2)
  \quad\Longrightarrow\quad
  \mathcal{R}(x_1) = \mathcal{R}(x_2).
\]
\end{definition}

The condition demands that whenever the projection declares two states
equivalent, those states must possess identical reachability structure. If
the projection collapses two states whose future horizons differ, then it has
destroyed a distinction that has genuine bearing on what the system can
subsequently do or become. Such a projection is inadmissible.

Several features of this definition deserve careful attention.

First, the criterion is future-directed rather than present-directed. Most
conventional criteria for evaluating abstractions assess them according to
their relationship to present observations or current performance metrics.
Admissibility evaluates whether the abstraction preserves the structure that
determines what futures remain available. This relocation from present to
future is not a superficial change of framing; it corresponds to a genuine
difference in what gets measured and what failures become visible.

Second, admissibility is intentionally demanding. It does not ask whether
the collapsed states are difficult to distinguish empirically, nor whether
the distinction improves short-horizon prediction, nor whether the collapse
is unlikely to matter in practice. It asks only whether the collapse alters
the geometry of future possibility. This demanding character is appropriate
because the failures that admissibility is designed to detect are precisely
the ones that conventional metrics miss: inadmissible collapses look
acceptable from the perspective of current performance and reveal their
consequences only when the futures they have obscured become relevant.

Third, admissibility furnishes a way of understanding why many successful
abstractions remain trustworthy despite discarding vast quantities of detail.
Fluid mechanics ignores the precise positions and velocities of individual
molecules. Thermodynamics ignores microscopic trajectories entirely.
Population genetics ignores individual organisms in favor of allele
frequencies. These abstractions succeed not despite their information losses
but because those information losses are, to a very good approximation,
admissible: the distinctions they collapse do not significantly alter the
reachability structure relevant to the macroscopic phenomena under
investigation.

\begin{proposition}
Let $\pi$ be an admissible projection. Then for any two states
$x_1, x_2 \in \mathcal{X}$ with $\pi(x_1) = \pi(x_2)$, any policy
$\sigma$ defined over $\mathcal{M}$ induces the same distribution over
future trajectories from $x_1$ as from $x_2$, in the sense that the
reachable sets under $\sigma$ are identical.
\end{proposition}

\begin{proof}
Suppose $\pi(x_1) = \pi(x_2)$ and $\pi$ is admissible. Then by definition
$\mathcal{R}(x_1) = \mathcal{R}(x_2)$. Any policy $\sigma : \mathcal{M}
\rightarrow \mathcal{A}$ assigns actions based solely on the representation,
and since $\pi(x_1) = \pi(x_2)$, the same action sequence is induced from
both states. The resulting trajectories draw upon the same reachable set.
Hence no outcome reachable from $x_1$ under $\sigma$ is unreachable from
$x_2$ under $\sigma$, and conversely.
\end{proof}

The proposition makes precise why admissibility matters operationally. A
policy operating in the representational space $\mathcal{M}$ cannot
distinguish states that $\pi$ has identified as equivalent. If those states
differ in their reachable futures, the policy will inevitably sometimes
attempt to reach a future that is not available from the actual state it
inhabits, or fail to avoid a future that the actual state makes inescapable.

A useful way to understand the criterion is to consider two types of
representational error. In the first type, the projection merges two states
whose observable properties differ slightly but whose reachable futures are
identical. Such a projection loses information about the present but does not
alter the geometry of future possibility. The collapse is imperfect in an
informational sense but remains admissible: nothing about what the system can
subsequently do has been obscured. In the second type, the projection merges
two states whose present observations are nearly identical but whose reachable
futures diverge dramatically. The crack in the bridge is the canonical
example: the observable difference may be tiny, the informational loss may
be negligible, and the effect on short-horizon prediction may be
imperceptible, yet the future consequences are enormous. This second collapse
is inadmissible.

Admissibility therefore reframes the traditional relationship between
abstraction and fidelity in an important way. A representation need not be
maximally detailed to be good, and a highly detailed representation is not
automatically trustworthy. What matters is not the quantity of information
preserved but the future structure preserved. The criterion is neither
completeness nor accuracy in the ordinary sense. It is the preservation of
reachability structure: the maintenance of distinctions whose loss would alter
what the system can subsequently do or become.

%%% SECTION 7
\section{Compression as Controlled Forgetting}

The concept of admissibility shifts attention from information preservation
to future preservation. This shift becomes especially important when
considering one of the most influential ideas in modern theories of
intelligence: the identification of understanding with compression. The
appeal of this identification is powerful and its intellectual heritage is
deep. Scientific laws compress observations into compact principles.
Mathematical equations replace vast collections of empirical facts with
concise descriptions. Statistical models reduce complexity by identifying
regularities. Machine learning systems discover latent representations that
permit more efficient prediction and storage. Across all of these domains,
successful understanding appears inseparable from successful compression, and
the intuition that an elegant, parsimonious description reveals something
deeper than a verbose one has been central to scientific and mathematical
practice for centuries.

The intuition has a rigorous formal tradition. Minimum description length
provides a Bayesian reformulation of Occam's razor in terms of code lengths.
Kolmogorov complexity defines the intrinsic information content of a string
as the length of the shortest program that generates it. Algorithmic
information theory unifies these notions and connects them to computational
universality. Building upon these foundations, compression-based theories of
intelligence interpret curiosity as the search for compressibility and
understanding as the discovery of shorter descriptions of experience.

The success of this viewpoint should not be underestimated. Many of the
greatest achievements of science can be understood as acts of compression,
and the insight that regularity and compressibility are closely related is
both genuine and important. The difficulty is not that compression is
unimportant but that it conceals a cost that is rarely examined as carefully
as the benefits.

Every successful compression scheme operates by forgetting.

A compression algorithm reduces description length by identifying distinctions
that need not be preserved. The bits that are eliminated correspond to
distinctions that the compression scheme has decided are redundant or
irrelevant for the purposes of reconstruction. Every abstraction, every
category, every scientific law, and every learned representation achieves
efficiency by discarding distinctions judged unnecessary for some purpose.
The success of compression is therefore inseparable from selective amnesia,
and the question of what has been forgotten is as important as the question
of what has been retained.

The conventional discussion focuses on what has been preserved. A compressed
representation is praised because it captures structure, reveals regularities,
or identifies hidden patterns. Less attention is devoted to the inverse
question: what is no longer distinguishable? From the perspective developed
in this essay, that question is fundamental.

Compression is not merely the preservation of structure. It is the destruction
of distinctions. Every reduction in description length corresponds to a
collapse of states that were previously represented separately. The critical
issue is therefore not whether compression occurs but which distinctions
disappear under the compression scheme.

This is where admissibility becomes operative. Traditional measures of
compression evaluate reductions in information, description length, or
predictive redundancy. Admissibility evaluates whether the distinctions
removed by compression alter the future accessibility structure of the
system. The difference is subtle but profound, and it becomes dramatic in
precisely the cases where it matters most.

A compression scheme may achieve extraordinary efficiency while remaining
admissible. This occurs when the distinctions it removes do not alter
reachable futures. Many scientific abstractions succeed precisely because they
discard details that are irrelevant to the future structure under
consideration: molecular positions can be forgotten in thermodynamics because
the macroscopic reachability structure is insensitive to them. However,
compression and admissibility are not equivalent, and there is no guarantee
that compressive efficiency and reachability preservation point in the same
direction.

A representation may become more compressive by eliminating distinctions that
appear insignificant from the perspective of present observations. Yet those
same distinctions may occupy critical positions within the geometry of future
possibility. The crack in the bridge is again the canonical example.
From the perspective of compression, the crack is an attractive candidate for
elimination: its contribution to present observations is tiny, its influence
on average behavior is negligible, and across a sufficiently large training
distribution, preserving the distinction between cracked and uncracked states
may appear wasteful. A compressed representation can often achieve lower
description length by treating them as effectively identical. The resulting
representation becomes simpler, and by every internal metric of the
compression scheme it has succeeded. It has also become blind to a
distinction that determines which futures are accessible.

The significance of the crack therefore reveals a broader principle: the
importance of a distinction cannot be determined solely by its contribution
to present description, nor by its contribution to average prediction, nor
by any local measure of its current influence. A distinction may appear
insignificant according to every local metric while remaining indispensable
to the preservation of future accessibility. This observation exposes a
limitation in compression-centered theories of intelligence that is not
merely practical but principled.

Suppose a representation achieves exceptional compression performance,
discovers elegant latent structure, minimizes redundancy, and reveals deep
regularities. None of these achievements guarantees that the representation
preserves the distinctions necessary for future accessibility, because
compression rewards forgetting while admissibility constrains it. The two
objectives frequently align, and when they do compression is not only
efficient but trustworthy. But they are not identical, and the pressure
exerted by aggressive compression is toward the elimination of distinctions
whose contribution to description length is small---which is precisely the
pressure that may eliminate reachability-critical features.

This permits a reformulation of the relationship between knowledge and
compression. Rather than asking how much a representation can forget while
retaining predictive power, one must ask how much it can forget while
retaining the future structure of the system. The question is therefore not
merely:
\begin{quote}
  How can the world be compressed?
\end{quote}
but rather:
\begin{quote}
  What cannot be forgotten without destroying the future?
\end{quote}

%%% SECTION 8
\section{Admissibility Distortion}

The concept of admissibility provides a binary criterion: a projection either
preserves reachability structure or it does not. Yet binary criteria, while
valuable for establishing the presence or absence of a property, are often
insufficient for practical analysis and theoretical development. Real
representations are neither perfectly admissible nor completely arbitrary.
They occupy an intermediate position in which some identifications are
harmless, others are consequential, and the degree of harm varies
continuously across the state space. A useful theory therefore requires not
only a criterion for perfect admissibility but also a measure of degree of
departure from it---a way of quantifying how much reachability structure a
given projection destroys.

The need for such a measure becomes evident when considering the limitations
of existing representational metrics. Information-theoretic measures evaluate
how much information remains after a projection. Statistical measures
evaluate explanatory power or predictive performance. Compression measures
evaluate reductions in description length. Optimization measures evaluate
success relative to some objective function. Each of these quantities
captures something important. None directly measures what futures have
disappeared, because the disappeared futures are precisely those that have
been rendered invisible by the projection and therefore cannot be assessed
from within the representational space.

To formalize a measure of destructive forgetting, consider a projection
$\pi : \mathcal{X} \rightarrow \mathcal{M}$ and let $\mathcal{R}(x)$
denote the set of futures reachable from state $x$. Let $d_H$ denote a
suitable distance measure between reachable future sets---for concreteness,
one may take $d_H$ to be a Hausdorff-type metric on the relevant space of
future trajectories, though the definition is compatible with optimal
transport distances and other metrizations of set-valued maps.

\begin{definition}[Admissibility Distortion]
The admissibility distortion of a projection $\pi$ is
\[
  D_A(\pi)
  \;=\;
  \mathbb{E}_{x_1,\,x_2}
  \!\left[
    \mathbf{1}_{\{\pi(x_1)=\pi(x_2)\}}
    \cdot
    d_H\!\left(\mathcal{R}(x_1),\,\mathcal{R}(x_2)\right)
  \right],
\]
where the expectation is taken over pairs $(x_1, x_2)$ drawn from an
appropriate distribution over $\mathcal{X} \times \mathcal{X}$.
\end{definition}

Conceptually, admissibility distortion measures the expected divergence
between future horizons for pairs of states that the projection has identified
as equivalent. The indicator function restricts attention to states that have
been collapsed by the projection. The distance term evaluates how different
their future horizons actually are, despite the representational merger.
The resulting quantity therefore measures not information loss in the abstract
but destructive forgetting specifically: the extent to which the projection
has collapsed distinctions that separate genuinely different futures.

Several properties of this definition clarify its relationship to existing
measures and its significance for the arguments developed elsewhere in this
essay.

\begin{proposition}
A projection $\pi$ is admissible if and only if $D_A(\pi) = 0$.
\end{proposition}

\begin{proof}
If $\pi$ is admissible, then $\pi(x_1) = \pi(x_2)$ implies
$\mathcal{R}(x_1) = \mathcal{R}(x_2)$, so $d_H(\mathcal{R}(x_1),
\mathcal{R}(x_2)) = 0$ for every pair satisfying the indicator. Hence
$D_A(\pi) = 0$. Conversely, if $D_A(\pi) = 0$, then for almost every
pair with $\pi(x_1) = \pi(x_2)$ we have $d_H(\mathcal{R}(x_1),
\mathcal{R}(x_2)) = 0$, which implies $\mathcal{R}(x_1) =
\mathcal{R}(x_2)$, so $\pi$ is admissible.
\end{proof}

The measure therefore extends the binary admissibility criterion continuously,
with zero distortion corresponding to perfect admissibility and larger values
corresponding to greater degrees of destructive forgetting.

A critical distinction separates admissibility distortion from conventional
information-theoretic measures of representational quality. Prediction error
asks whether future observations can be forecast accurately from the
representation. Compression measures ask whether observations can be
described economically within the representation. Information-theoretic
measures ask whether distinctions remain statistically resolvable. All of
these quantities evaluate properties relative to the representation and its
relationship to observed data. Admissibility distortion asks a structurally
different question:

\begin{quote}
  How much future accessibility has been destroyed by the act of
  abstraction itself?
\end{quote}

This distinction reveals why apparently successful representations may fail
in ways that remain entirely invisible to conventional evaluation procedures.
A model may achieve excellent predictive performance, state-of-the-art
compression, and robust downstream task success while exhibiting significant
admissibility distortion. Such a model appears successful because its
failures occur not within the representational space---where evaluation
takes place---but at the boundaries between futures that the representation
has rendered indistinguishable.

The crack in the bridge again provides the illustrative case. A
representation that identifies two bridge states as equivalent because their
observable behavior is nearly identical throughout the available training data
will exhibit low prediction error and good compression. Yet if one state
contains a microscopic structural defect and the other does not, the
reachable futures may differ enormously, and $D_A(\pi)$ will reflect this
even when all ordinary performance metrics look favorable.

Admissibility distortion also clarifies the relationship between abstraction
and structural risk. Every abstraction introduces the possibility of
distortion, and the question is not whether distortion exists but where it
is concentrated. In many systems the greatest distortion occurs near critical
boundaries: regions where small differences produce large changes in future
accessibility. These boundaries include structural failure thresholds,
ecological tipping points, software vulnerabilities, financial crises,
evolutionary bifurcations, and many other phenomena in which future
possibility changes discontinuously. From the perspective of reachability,
such boundaries are precisely where inadmissible abstraction is most
dangerous, and they tend to be the regions where statistical sparsity makes
conventional evaluation metrics most misleading.

The following result shows that admissibility distortion is, in a precise
sense, monotone under the composition of projections. Additional layers of
abstraction cannot reduce distortion that has already been introduced.

\begin{theorem}[Distortion Monotonicity Under Composition]
Let $\pi_1 : \mathcal{X} \to \mathcal{M}$ and $\pi_2 : \mathcal{M} \to
\mathcal{N}$ be projections and let $\pi = \pi_2 \circ \pi_1 :
\mathcal{X} \to \mathcal{N}$. Under mild regularity conditions on the
sampling distribution,
\[
  D_A(\pi_2 \circ \pi_1) \;\geq\; D_A(\pi_1).
\]
\end{theorem}

\begin{proof}
Any pair $(x_1, x_2)$ with $\pi_1(x_1) = \pi_1(x_2)$ necessarily satisfies
$\pi_2(\pi_1(x_1)) = \pi_2(\pi_1(x_2))$, so the indicator
$\mathbf{1}_{\{\pi(x_1)=\pi(x_2)\}}$ takes value 1 whenever
$\mathbf{1}_{\{\pi_1(x_1)=\pi_1(x_2)\}}$ does. The expectation defining
$D_A(\pi_2 \circ \pi_1)$ therefore integrates over a superset of the pairs
contributing to $D_A(\pi_1)$, and each contributing pair registers the same
reachability divergence $d_{\mathcal{R}}(\mathcal{R}(x_1),
\mathcal{R}(x_2))$. Since all terms are nonnegative, the composed distortion
is at least as large as the initial distortion.
\end{proof}

The practical implication is significant: once an inadmissible collapse has
been introduced at the level of $\pi_1$, no downstream representational
transformation $\pi_2$ can undo it. Every further layer of abstraction
applied to an already-inadmissible representation either preserves or
increases the distortion. Remediation must therefore occur upstream, at the
level of the original projection, rather than downstream through refinement
of the model built upon it.

%%% SECTION 9
\section{The Proxy Problem}

The preceding sections have developed the positive core of the reachability
framework: the ontological priority of accessibility, the formal concept of
admissibility, the measure of admissibility distortion, and the contrast
between compression as information reduction and compression as controlled
forgetting. The present section applies this framework to a cluster of
related failures that share a common structure.

Each of the dominant modern frameworks for evaluating representations and
intelligence---compression, prediction, statistical significance, and
optimization---defines a proxy metric and then, through the momentum of
successful application, allows that metric to migrate from instrument to
criterion. The proxy is mistaken for the thing it was originally designed to
track. Once this migration has occurred, the failure mode is always the same:
the system performs excellently according to its own metric while destroying
precisely the structure the metric was originally intended to preserve.

The formal structure of the problem can be stated as a general
non-equivalence. Let $M(x)$ denote any of the proxy metrics under
consideration (compression rate, prediction accuracy, statistical
significance, expected reward), and let $A(x)$ denote the admissibility
of the representation with respect to state $x$. Then in general
\[
  M(x_1) = M(x_2)
  \;\not\Rightarrow\;
  \mathcal{R}(x_1) = \mathcal{R}(x_2),
\]
which is to say: equality of proxy metric values does not entail equality of
reachable future sets. This non-implication is the single theorem that the
following subsections instantiate across different domains.

\subsection{Prediction as a Derived Activity}

The modern history of artificial intelligence has elevated prediction to a
position of unusual prominence. Predictive accuracy serves as the dominant
metric in machine learning, scientific modeling, and increasingly in theories
of intelligence itself. Systems are trained to predict the next token, the
next frame, the next state, or the next action. Success is measured by the
extent to which future observations can be anticipated from present ones. The
practical achievements of this paradigm are undeniable: predictive systems
have transformed language modeling, computer vision, scientific simulation,
weather forecasting, protein structure prediction, and many other domains.

Yet predictive success and representational adequacy are not identical. The
reason is that prediction evaluates a model according to its performance
within a representation, not according to the admissibility of the
representation itself. Prediction assumes a partition of the world and then
asks how accurately dynamics can be modeled inside that partition. The
partition itself remains largely outside the scope of evaluation.

Many contemporary frameworks implicitly adopt the inference chain
\[
  \text{Good Representation}
  \;\Rightarrow\;
  \text{Good Prediction}
  \;\Rightarrow\;
  \text{Good Planning}
  \;\Rightarrow\;
  \text{Good Control}.
\]
The plausibility of the chain derives from the practical success of
predictive learning. If a representation permits accurate forecasting, it
appears natural to assume that the representation has captured what matters.
Yet the non-equivalence established above shows why this inference is
unwarranted. Prediction concerns the evolution of distinctions that remain
available within the representation. It says nothing about distinctions that
have already been collapsed.

Let $P(x)$ denote the predictive equivalence class associated with state
$x$: two states are predictively equivalent if they produce essentially the
same predictions according to the model. Then predictive equivalence does not
imply reachability equivalence:
\[
  P(x_1) = P(x_2)
  \;\not\Rightarrow\;
  \mathcal{R}(x_1) = \mathcal{R}(x_2).
\]
There may exist states $x_1, x_2$ that appear identical from the perspective
of prediction yet occupy different positions within the geometry of future
possibility. The significance of this gap is easy to underestimate because
predictive performance remains excellent until the moment the distinction
becomes important. The crack propagates. The software fault activates. The
ecological threshold is exceeded. The predictive system has not erred in any
ordinary sense. It has inherited a representation that treated the critical
distinction as irrelevant.

\subsection{Statistics as Local Importance}

Modern science relies heavily upon statistical measures of importance:
variance explained, effect size, correlation strength, information gain, and
related quantities. Such measures have proven indispensable for distinguishing
signal from noise in complex systems. Yet the success of statistical methods
has encouraged a subtle assumption that importance and frequency are
fundamentally aligned.

The assumption is false in general. Statistical significance concerns the
contribution of a feature to observed variation. Admissibility concerns the
contribution of a feature to future accessibility. The two quantities are
often correlated. They are not identical.

The distinction becomes visible whenever rare events reshape the future. A
mutation occurring in a single organism may eventually transform an entire
population. A software fault may remain dormant across millions of executions
before triggering catastrophic failure. Statistical frameworks tend to
discount such phenomena because their contribution to aggregate variation is
limited. Their significance derives not from frequency but from their location
within the geometry of future possibility.

The asymmetry can be stated formally. Statistical importance is often
proportional to a quantity of the form
\[
  I_S(x) \;\propto\; \Pr(x) \cdot E(x),
\]
where $\Pr(x)$ denotes frequency and $E(x)$ denotes a measure of effect
within the current distribution. Admissibility importance depends upon
\[
  I_A(x) \;\propto\; \Delta \mathcal{R}(x),
\]
the change in reachable futures induced by the state. These quantities are
not generally equivalent. A state may possess low statistical importance
while exerting enormous influence on future accessibility, and the disparity
is greatest near thresholds and singular transitions---precisely the regions
where the future changes most discontinuously.

This observation extends to compression-based theories of curiosity.
Curiosity oriented primarily toward compressibility will allocate attention
to regions where better regularities are found and description length can be
reduced. Yet the most important discoveries need not be those that maximize
compressibility. A crack in a bridge is not interesting because it improves
compression; a pathogen is not significant because it reveals a new
regularity. They matter because they expose regions where future
accessibility changes abruptly. An admissibility-oriented intelligence would
allocate attention differently, seeking not merely regions where prediction
improves but regions where the structure of future possibility changes.

\subsection{Optimization Within Reachability Geometry}

Among the intellectual frameworks that dominate contemporary thought, few
enjoy the prestige of optimization. Economics is built around it. Control
theory is built around it. Modern machine learning has elevated optimization
to an almost universal explanatory principle. The success of this approach
is understandable: given a space of possible states and a criterion for
evaluating them, optimization offers a systematic procedure for selecting
preferred outcomes. Whether the objective is profit, efficiency, accuracy,
fitness, or expected reward, the underlying mathematical structure is
remarkably unified.

Yet optimization possesses a hidden dependency that is frequently overlooked.
Optimization does not create possibilities. It searches among them. An
optimizer can choose among available futures, rank trajectories by objective
value, and exploit landscape structure. It cannot determine what futures
exist, what trajectories are available, or what landscape structure is
present. These are given by the geometry of the system, and that geometry
is the product of the representation.

Suppose an optimizer is given a representation that has collapsed a
reachability-critical distinction. The optimizer then performs its search
within a geometry that no longer contains the distinction. It may optimize
the objective perfectly within that geometry. It will also amplify the
representational error, because optimization is precisely the activity of
finding extreme solutions---solutions near the boundaries of the feasible
set, where small geometric inaccuracies become large operational ones.

This observation suggests a reinterpretation of familiar optimization
failures. The standard explanation attributes failure to misspecified
objectives. A deeper class of failure exists in which the objective is
entirely reasonable but the representation is inadmissible. In such cases
the optimizer faithfully pursues its objective within a geometry that has
already forgotten something essential. The resulting behavior may satisfy
every local criterion while systematically destroying options that the
representation has rendered invisible. The optimizer does not malfunction;
the map is wrong.

\subsection{Representation Without Accessibility}

The ideal of the digital twin represents the limiting case of
representational thinking. Construct a sufficiently accurate representation
of a system, update it continuously with incoming data, and one obtains a
virtual counterpart capable of prediction, diagnosis, optimization, and
control. The digital twin aspires to be so faithful that the distinction
between representation and reality begins to dissolve.

Yet representational fidelity and admissibility are different properties, and
the pursuit of fidelity does not guarantee the preservation of reachability
structure. Fidelity is fundamentally a descriptive notion: it concerns
correspondence between model and system, evaluated by how closely the model's
outputs match the system's observed behavior. Admissibility concerns something
different: the preservation of distinctions that determine what futures remain
possible.

Consider two representations of the same bridge. The first contains
exhaustive information about present conditions: stress distributions,
material composition, temperature gradients, traffic patterns, and structural
geometry with remarkable precision. The second contains far less information,
but tracks with particular care the boundaries where future accessibility
changes---locations where cracks propagate, thresholds where structural
integrity fails, and transitions where repair becomes impossible.

From the perspective of fidelity, the first model appears preferable. From
the perspective of admissibility, the answer is less obvious. If the first
model omits the distinctions that determine accessibility boundaries while
the second preserves them, then the simpler model may provide more meaningful
guidance about the future structure of the system, even though it is less
faithful in the conventional sense.

The dream of the perfect digital twin also conceals a deeper impossibility.
Completeness is not merely impractical; it is conceptually incoherent.
A model that preserved every distinction would cease to function as a model.
It would become indistinguishable from the system itself. Representation
requires collapse. The question is always which collapses occur, and the
digital twin paradigm answers this question by maximizing fidelity when it
should be asking about admissibility.

\subsection{The Accessibility Estimation Problem}

The framework developed in this section raises a practical question that
deserves acknowledgment rather than evasion: if admissibility distortion is
what matters, how can it be estimated without direct access to
$\mathcal{R}_D(x)$? In many real systems the full reachability structure is
not directly observable and cannot be computed in closed form.

This is not a fatal difficulty for the framework; it is an open research
problem that the framework makes precise. Several approaches are available
at different levels of fidelity. Intervention experiments and controlled
perturbation studies probe the boundary between accessible and inaccessible
futures by actively perturbing states and observing whether trajectories
diverge. Control-theoretic reachability analysis, viability kernels, and
barrier certificate methods provide formal inner and outer approximations to
$\mathcal{R}_D(x)$ for systems with known dynamics. Counterfactual simulation
and model-based rollout estimate reachable regions by sampling trajectories
under learned or specified dynamics. Empowerment---the mutual information
between actions and reachable states---provides an information-theoretic
proxy that is computable from transition models or from data. Basin-of-attraction
estimation and tipping-point detection offer reachability-adjacent quantities
that can be estimated from time-series data near critical thresholds.

None of these approaches provides exact admissibility distortion, and the
problem of estimating $D_A(\pi)$ without full knowledge of $\mathcal{R}_D$
is genuinely hard. But the difficulty of computation does not undermine the
conceptual primacy of the criterion any more than the difficulty of computing
optimal strategies undermines game theory, or the difficulty of measuring
entropy production undermines thermodynamics. The framework specifies what
should be preserved; the engineering challenge is to construct tractable
proxies that track it. Framing the problem clearly is the prerequisite for
solving it.

%%% SECTION 10
\section{Intelligence as Reachability Preservation and Navigation}

The preceding sections have pursued a largely critical task, showing that
compression, prediction, statistical significance, optimization, and
representational fidelity are each insufficient as foundational criteria for
intelligence or representational adequacy. The recurring pattern is now
visible. Every framework evaluates success through a proxy metric.
Compression rewards shorter descriptions. Prediction rewards accurate
forecasts. Statistics rewards explanatory power. Optimization rewards
objective attainment. Fidelity rewards resemblance. These metrics often
correlate with successful interaction with the world, which explains their
practical value.

The mistake occurs when correlation is elevated into ontology. The proxies
are treated as foundations rather than as instruments whose validity depends
upon a prior condition. The result is a persistent tendency to confuse the
instrument with the structure it was designed to track, and to optimize the
instrument while allowing the structure to degrade invisibly.

The reachability framework suggests a different starting point. If the
significance of a state derives from the futures it permits, and if the
adequacy of a representation derives from the futures it preserves, then the
central question of intelligence changes. Intelligence can no longer be
defined primarily as prediction, compression, optimization, or
representation. Each of those activities becomes subordinate to a more
fundamental concern: the preservation of future accessibility.

This proposal may initially appear unusual because intelligence is
traditionally described in epistemic terms. An intelligent system is assumed
to possess knowledge, discover regularities, predict outcomes, solve
problems, and reason effectively about its environment. All of these
descriptions contain an important truth. Yet they share a common limitation:
they focus on internal competence rather than on the relationship between
the system and the future structure of its environment.

An intelligent system is not merely one that knows many things, for a
perfectly well-informed system that systematically drives itself toward
irreversible dead ends would be difficult to call intelligent in any
meaningful sense. Nor is it merely a system that predicts accurately, for
a system with exceptional short-horizon predictive capability might still
fail to preserve the distinctions necessary for long-term adaptability.
Nor is it merely an efficient compressor, for compression that eliminates
reachability-critical distinctions may be representationally elegant while
being operationally catastrophic.

The inadequacy of these descriptions becomes especially apparent when
considering biological organisms. Many organisms survive for astonishingly
long periods under conditions of extreme uncertainty, with limited predictive
capabilities, crude internal models, incomplete sensory information, and
severe cognitive and energetic constraints. The reason is not perfect
prediction. The reason is that successful behavioral strategies tend,
robustly and across a wide range of conditions, to preserve future
accessibility. Successful organisms avoid irreversible traps. They maintain
flexibility. They retain multiple behavioral options. They protect the
conditions under which future adaptation remains possible. They do not
necessarily know what will happen; they preserve the ability to respond
constructively when it does.

This observation suggests a reinterpretation of intelligence. Intelligence
is not fundamentally the capacity to predict the future. It is the capacity
to remain capable of acting within it. The distinction is subtle but
important. Prediction concerns accuracy relative to an expected future.
Reachability concerns the maintenance of possibility across a range of
futures that may or may not correspond to expectations.

The same reinterpretation applies to the subordinate cognitive activities.
Learning, on the reachability account, is the refinement of
admissibility-preserving distinctions: a system learns when it improves its
ability to recognize the boundaries separating different future horizons.
Memory functions as a mechanism for preserving distinctions across time whose
loss would reduce future adaptability. Reasoning is navigation through spaces
of possibility, with the purpose of evaluating trajectories, identifying
constraints, and preserving access to desirable futures rather than merely
deriving conclusions. Safety is the preservation of accessibility under
uncertainty, distinguishing dangerous states as those that collapse future
possibility.

What unites these activities is not information in the abstract but future
accessibility, and the formal condition that allows intelligence to preserve
future accessibility through abstraction is admissibility. A representation
is useful because it preserves the distinctions necessary for navigation
within the reachability structure of the world. Without admissibility,
prediction is blind to the futures that matter most. Without admissibility,
optimization amplifies the errors introduced by inadmissible collapse.
Without admissibility, compression destroys exactly the distinctions that
would be needed when the compressed states diverge into different futures.

The positive thesis of this essay can therefore be stated precisely.

\begin{quote}
  The purpose of intelligence is not to model the world as accurately as
  possible. The purpose of intelligence is to preserve, discover, and
  navigate the futures that matter, through abstraction that remains
  admissible with respect to the reachability structure of the environment.
\end{quote}

The word \emph{navigate} is deliberate and important. A purely preservationist
account of intelligence would leave it indistinguishable from mere
inertia---a system that never acts also never destroys future options. What
distinguishes intelligence from passive persistence is the capacity to move
through the reachability structure in ways that maintain or expand the space
of available futures while avoiding irreversible collapse. Preservation is
the enabling condition; navigation is the activity; and admissibility is the
constraint that keeps navigation from inadvertently destroying the geometry
it depends upon.

%%% SECTION 11
\section{Abstraction After the Inversion}

At this point a potential misunderstanding must be addressed directly. The
argument developed throughout this essay may appear to constitute an attack
upon abstraction itself. Compression has been criticized. Prediction has been
criticized. Statistical significance has been criticized. Optimization has
been criticized. Representational fidelity has been criticized. A reader
might reasonably conclude that the reachability framework demands the
preservation of every distinction and therefore rejects abstraction
altogether.

Such a conclusion would be mistaken, and it would be the opposite of what
the framework intends.

The reachability thesis is not an argument against abstraction. It is an
argument about the conditions under which abstraction succeeds. No finite
intelligence can operate without abstraction. Every organism, every
scientific theory, every engineering discipline, every mathematical model,
and every artificial intelligence system depends upon the ability to ignore
vast quantities of detail. The world contains more distinctions than any
finite system can represent explicitly, and the attempt to preserve all of
them would not produce a better representation but simply no representation
at all. Abstraction is not optional. It is the price of cognition, and it
is paid with necessity by any system that must act within a complex world.

The central question has never been whether distinctions should be removed.
The central question is which distinctions may be removed safely---that is,
without altering the future structure that the system must navigate. This
is precisely why admissibility occupies a fundamental role in the framework:
an admissible abstraction is not one that preserves all information but one
that preserves the future structure relevant to the domain under
consideration.

Seen from this perspective, many of the greatest achievements of science
become examples of successful admissible projection. Fluid mechanics provides
a particularly instructive case. A fluid description forgets almost
everything about individual molecules: the precise position and velocity of
each particle disappears entirely, and the overwhelming majority of
microscopic distinctions are collapsed into aggregate quantities such as
density, pressure, and velocity fields. By ordinary informational standards,
the loss is immense. Yet fluid mechanics remains extraordinarily successful
as a predictive and explanatory framework, not despite the loss but because
of it. The distinctions that fluid mechanics collapses do not significantly
alter the macroscopic reachability structure under investigation. The
projection is admissible at the macroscopic level of description.

The same pattern appears throughout the history of successful scientific
abstraction. Thermodynamics forgets microscopic trajectories while preserving
the future structure of macroscopic state variables. Population genetics
forgets individual organisms while preserving the dynamics of allele
frequencies at the population level. Celestial mechanics forgets the internal
molecular structure of planets while preserving the accessibility structure
of orbital motion. Classical electrodynamics forgets the quantum mechanical
details of charge distributions while preserving the macroscopic field
structure. Each framework succeeds because its abstractions are admissible
with respect to the domain they are intended to describe. They are not
arbitrary compressions; they are projections that preserve the future
geometry at the relevant level of description.

This observation clarifies a recurring confusion in contemporary discussions
of artificial intelligence. Representational systems are often evaluated
according to how much information they retain, how accurately they predict,
or how effectively they support downstream tasks. These metrics implicitly
assume that preserving information is the primary challenge and that useful
behavior emerges naturally once sufficient information remains available.
The reachability framework reverses this logic. Information becomes valuable
because certain distinctions support the preservation of future
accessibility. A representation is useful not because it contains information
in the abstract but because it preserves distinctions that matter for the
future structure of the system it represents.

This inversion changes the meaning of abstraction itself. Abstraction is no
longer understood as the removal of detail for the sake of efficiency. It
becomes the selective preservation of future structure. A successful
abstraction is therefore not one that forgets little. It is one that forgets
wisely. The traditional evaluation question asks how much information can be
removed while retaining acceptable performance. The admissibility question
asks which distinctions can be removed without altering what futures remain
reachable. The two inquiries overlap, but they are not equivalent, and the
cases where they diverge are precisely the cases where the most important
failures occur.

%%% SECTION 12
\section{The Hierarchy Reversed}

The various arguments of this essay, pursued through the specific domains of
compression, prediction, statistics, optimization, and representation, can
now be unified into a single structural observation. Every framework that
has been criticized exhibits the same fundamental ordering error. It places
representation, information, or prediction at the foundation of the
conceptual hierarchy and treats reachability as a downstream consequence.
The reachability thesis reverses this ordering.

The dominant intellectual picture of the modern era places information and
representation near the foundation of inquiry. A system acquires information
about the world. That information supports the construction of
representations. Representations support prediction. Prediction supports
planning. Planning supports control. Schematically:
\[
  \text{Information}
  \;\rightarrow\;
  \text{Representation}
  \;\rightarrow\;
  \text{Prediction}
  \;\rightarrow\;
  \text{Planning}
  \;\rightarrow\;
  \text{Control}.
\]
Within this framework, the central challenge becomes the improvement of
representations. Better representations yield better predictions. Better
predictions yield better plans. Better plans yield better control. Progress
is understood primarily as progress in representation and prediction.

The difficulty is that every stage of this hierarchy presupposes something
that the hierarchy itself does not explain. The representation must already
preserve the distinctions that matter. The prediction must already occur
within a meaningful space of futures. The planner must already possess
access to an appropriate geometry of possibilities. The controller must
already operate within a representation whose distinctions remain relevant
to future accessibility. In every case the framework assumes the preservation
of future structure without making that preservation an explicit object of
analysis.

The reachability perspective therefore proposes a different ordering:
\[
  \text{Reachability}
  \;\rightarrow\;
  \text{Admissibility}
  \;\rightarrow\;
  \text{Representation}
  \;\rightarrow\;
  \text{Prediction}
  \;\rightarrow\;
  \text{Planning}
  \;\rightarrow\;
  \text{Control}.
\]
Reachability occupies the foundation because it defines the future structure
of the system---the geometry within which every subsequent activity takes
place. Admissibility follows because it determines whether that future
structure survives the act of abstraction. Representation follows because
representations are projections whose validity depends upon admissibility.
Prediction, planning, and control follow as increasingly downstream
activities whose success depends upon every preceding level.

The reversal transforms the interpretation of every concept in the hierarchy.
Information ceases to be a primitive quantity and becomes valuable because
it supports admissible distinctions. The importance of a feature derives not
from the number of bits it encodes but from the role it plays in preserving
future accessibility. Representation ceases to be evaluated primarily
according to fidelity, compression, or predictive performance; these
quantities remain useful but become secondary indicators. The primary
question concerns admissibility: does the representation preserve the future
structure of the system?

Prediction likewise changes status. It remains valuable, but its value
becomes instrumental rather than foundational. Predictions assist navigation
within a reachability landscape; they do not define the landscape.
Optimization becomes perhaps the clearest example of the inversion.
Traditional accounts frequently treat optimization as the essence of
intelligence, but the reachability hierarchy reveals that optimization is
among the most downstream activities in the entire chain. Before optimization
can occur meaningfully, the future structure must exist; before that
structure can be represented, it must survive projection; before projection
can be evaluated, reachability must already be defined.

This reversal also clarifies why certain failure modes appear repeatedly in
modern technical systems. When a representation collapses a
reachability-critical distinction, optimization amplifies the error rather
than correcting it. When prediction ignores a future boundary, planning
inherits the blindness. When information is compressed inadmissibly, control
becomes fragile in precisely the regions where it matters most. Each failure
propagates downward through the hierarchy, and the root cause lies not at the
level of optimization or prediction but at the level of admissibility.

The hierarchy also provides a unified account of intelligence itself.
Learning becomes the discovery of admissible distinctions. Memory becomes
the preservation of admissible distinctions across time. Reasoning becomes
the evaluation of trajectories through admissible spaces of possibility.
Planning becomes the selection of trajectories within those spaces. Safety
becomes the preservation of accessibility under uncertainty and the avoidance
of irreversible collapse. The various cognitive and computational activities
traditionally treated as distinct are revealed as special cases of a common
process: maintaining contact with the future structure of the world through
abstractions that do not destroy the geometry they are intended to represent.

%%% SECTION 13
\section{What Cannot Be Forgotten?}

The argument of this essay began with a crack in a bridge.

The crack was small. It contributed little to observable behavior. It
explained almost none of the variance in the system under any standard
statistical framework. It added negligible predictive power. It compressed
poorly. A sufficiently aggressive abstraction could easily treat the cracked
bridge and the uncracked bridge as effectively identical, and by every
conventional metric such a treatment would appear justified. The crack was,
in the language of representational learning, a candidate for elimination.

Yet the crack mattered---and it mattered in a way that none of the
conventional metrics could register. It mattered because it occupied a
position within the geometry of future possibility such that the
futures accessible from the cracked state and the futures accessible from
the uncracked state diverged dramatically. The significance of the crack was
not determined by its influence upon present description but by its capacity
to alter what futures remained reachable. It separated futures. That is what
made it indispensable.

This simple observation has guided every stage of the argument, and its
implications are now fully visible.

The ontology of objects gives way to an ontology of reachability. Objects,
on the traditional picture, are primary and futures are derived from them.
On the reachability picture, futures are primary and objects emerge as
stabilized patterns within a deeper geometry of accessible transformations.
A bridge is not merely a collection of structural components; it is a
location within a network of possible futures, and its significance is
determined by which of those futures it permits and which it forecloses.

Every act of knowing becomes an act of identification. To understand
something is to collapse distinctions---to decide which features matter and
which may be treated as equivalent. Every ontology, every scientific law,
every representation, and every category is a theory of permissible
identification. The question is never whether collapse occurs but whether
the collapses that occur are admissible.

Abstraction becomes controlled collapse. Compression becomes controlled
forgetting. Both are necessary, both are inescapable, and both carry the
risk of inadmissibility---of collapsing distinctions whose loss
alters the geometry of future possibility.

Prediction becomes navigation within an already-existing geometry, and its
failures arise not when the navigation is performed poorly but when the
geometry upon which it operates has already lost essential structure.
Optimization becomes search inside an admissible geometry, and its most
dangerous failures occur not when the objective is wrong but when the
representational geometry within which the search occurs has collapsed
distinctions that determine which futures remain open.

Intelligence becomes the preservation of future accessibility rather than
the accumulation of predictive power or representational fidelity. It is
not a prediction engine, not a compression engine, and not an optimization
engine. It is a reachability-preservation engine: a system capable of
maintaining sensitivity to the distinctions that determine which futures
remain open and which futures become inaccessible, through abstractions
that do not destroy the structure they represent.

The recurring lesson is always the same. What matters most about a state is
not necessarily what contributes most strongly to present description. The
distinction may be statistically negligible. It may contribute almost nothing
to prediction. It may resist compression. It may occupy only a tiny region
of the representational space. Yet if it separates futures---if removing
it alters the geometry of accessible transformations---then its significance
cannot be measured by any quantity defined solely in the present.

This realization exposes a limitation that is not merely practical but
structural. Modern intellectual culture has become extraordinarily effective
at measuring what is frequent, predictable, compressible, and optimizable.
Entire fields are organized around these quantities. The critique offered
here is not a rejection of those achievements. It is a reminder that none of
those quantities is foundational. Each is valuable precisely because it is
often correlated with the preservation of future structure. Each fails
precisely when the correlation breaks down---which is to say, in the most
important cases: near the thresholds, at the tipping points, in the presence
of latent defects, when the crack propagates.

The problem is not that abstraction fails. The problem is that abstraction
always forgets, and the question of what it forgets has not been treated
with the seriousness it deserves. The framework developed here provides a
criterion for addressing that question. The value of a distinction is
determined not by its frequency, not by its contribution to prediction, not
by its compressibility, and not by its relationship to an objective function.
Its value is determined by the futures it preserves.

The question that can now be stated with precision is not merely
philosophical. Given a state space $\mathcal{X}$, a reachability structure
$\mathcal{R}$, and a projection $\pi : \mathcal{X} \rightarrow \mathcal{M}$,
which equivalence classes induced by $\pi$ satisfy the admissibility
condition, and how much admissibility distortion $D_A(\pi)$ is introduced
by those that do not? The measure of admissibility distortion provides a
formal handle on what is being destroyed, making it possible to evaluate
abstractions not only by what they reveal but by what futures they eliminate.

This transforms the philosophical question into a research program. The
scientist constructing a theory, the engineer building a model, the
statistician selecting variables, the machine learning system learning a
latent representation, the planner choosing actions, the optimizer searching
for solutions---each faces the same fundamental challenge, and the
reachability framework provides a common language in which that challenge
can be stated precisely. Each must decide what may be forgotten. The
framework supplies the criterion by which that decision can be evaluated.

The crack in the bridge is therefore more than an engineering example. It is
a metaphor for the entire argument, and it is also an instance of a general
structure that appears wherever a system must abstract from the world in
order to act within it. The most important features of a system are often
those that appear smallest from the perspective of the present. They are the
distinctions that occupy critical positions within the geometry of future
accessibility. They are the details whose removal transforms one future into
another. To forget them is not merely to lose information. It is to lose
worlds.

The deepest question of intelligence is therefore not epistemological and not
computational in the first instance. It is ontological. Before one asks what
can be predicted, compressed, optimized, represented, or controlled, one
must first establish what reachability structure exists, what abstractions
preserve it, and what cannot be forgotten without destroying the future.
That question---
\begin{quote}
  What cannot be forgotten without destroying the future?
\end{quote}
---is the question that has been haunting the argument since the first section.
It now has a formal home. It is the question of admissibility distortion,
and it is the central question of any theory of intelligence, knowledge, or
control that takes seriously the ontological priority of reachability.

\bigskip

\noindent\textbf{Acknowledgements.} The author thanks the many researchers,
theorists, engineers, and critics whose work, whether embraced or opposed,
helped clarify the distinction between prediction and possibility, between
compression and forgetting, and between the representation of the world and
the preservation of its futures. Any inadmissible collapses that remain in
the argument are entirely the author's own.


\appendix

%%% APPENDIX A
\section{Reachability, Projection, and Admissibility}

This appendix provides a more formal treatment of the reachability framework
used throughout the essay. Let $\mathcal{X}$ be a measurable state space and
let $\mathcal{U}$ be a class of admissible controls, interventions, policies,
or transformations. For each state $x \in \mathcal{X}$, define the reachable
future set
\[
  \mathcal{R}(x)
  =
  \{y \in \mathcal{X} : y \text{ is reachable from } x
  \text{ under some } u \in \mathcal{U}\}.
\]
The reachability relation induced by $\mathcal{U}$ is
\[
  x \leadsto y
  \quad\Longleftrightarrow\quad
  y \in \mathcal{R}(x).
\]
A representation is a measurable projection
\[
  \pi : \mathcal{X} \to \mathcal{M},
\]
where $\mathcal{M}$ is a representation space. The projection induces an
equivalence relation on $\mathcal{X}$:
\[
  x_1 \sim_\pi x_2
  \quad\Longleftrightarrow\quad
  \pi(x_1) = \pi(x_2).
\]
The projection is admissible exactly when this equivalence relation is finer
than reachability equivalence.

\begin{definition}[Reachability Equivalence]
Two states $x_1, x_2 \in \mathcal{X}$ are reachability-equivalent if
\[
  \mathcal{R}(x_1) = \mathcal{R}(x_2),
\]
written $x_1 \sim_{\mathcal{R}} x_2$.
\end{definition}

\begin{definition}[Admissible Projection---Formal]
A projection $\pi : \mathcal{X} \to \mathcal{M}$ is admissible if
\[
  x_1 \sim_\pi x_2
  \quad\Longrightarrow\quad
  x_1 \sim_{\mathcal{R}} x_2,
\]
equivalently, if $\pi(x_1) = \pi(x_2)$ implies $\mathcal{R}(x_1) =
\mathcal{R}(x_2)$.
\end{definition}

Thus admissibility requires that the projection collapse only those distinctions
that do not alter reachable futures.

\begin{proposition}[Quotient Characterization]
A projection $\pi : \mathcal{X} \to \mathcal{M}$ is admissible if and only if
there exists a well-defined map
\[
  \widehat{\mathcal{R}} : \pi(\mathcal{X}) \to \mathcal{P}(\mathcal{X})
\]
such that $\widehat{\mathcal{R}}(\pi(x)) = \mathcal{R}(x)$ for every
$x \in \mathcal{X}$.
\end{proposition}

\begin{proof}
Suppose first that $\pi$ is admissible. Define
$\widehat{\mathcal{R}}(m) = \mathcal{R}(x)$ for any $x$ such that
$\pi(x) = m$. If $\pi(x_1) = \pi(x_2) = m$, admissibility gives
$\mathcal{R}(x_1) = \mathcal{R}(x_2)$, so $\widehat{\mathcal{R}}$ is
independent of the chosen representative and hence well-defined. Conversely,
suppose such a map exists. If $\pi(x_1) = \pi(x_2)$, then
\[
  \mathcal{R}(x_1)
  = \widehat{\mathcal{R}}(\pi(x_1))
  = \widehat{\mathcal{R}}(\pi(x_2))
  = \mathcal{R}(x_2),
\]
so $\pi$ is admissible.
\end{proof}

This proposition gives a precise meaning to the central claim that an
admissible representation preserves future structure: reachability descends
to the quotient space induced by the representation if and only if the
representation is admissible.

%%% APPENDIX B
\section{Admissibility Distortion: Formal Development}

Perfect admissibility is often too strong a requirement for practical systems.
One therefore needs a quantitative measure of failure. Let $d_{\mathcal{R}}$
be a distance or divergence between reachable future sets. In many settings
this may be taken as the Hausdorff distance induced by a background metric
$d_{\mathcal{X}}$ on $\mathcal{X}$:
\[
  d_H(A, B)
  =
  \max
  \!\left\{
    \sup_{a \in A} \inf_{b \in B} d_{\mathcal{X}}(a, b),\;
    \sup_{b \in B} \inf_{a \in A} d_{\mathcal{X}}(a, b)
  \right\}.
\]
Alternative choices include optimal-transport distances between distributions
over future trajectories and task-dependent reachability divergences. The
definition of admissibility distortion is compatible with any of these choices.

\begin{definition}[Admissibility Distortion---General]
The admissibility distortion of a projection $\pi : \mathcal{X} \to \mathcal{M}$
is
\[
  D_A(\pi)
  =
  \mathbb{E}_{x_1,\,x_2}
  \!\left[
    \mathbf{1}_{\{\pi(x_1) = \pi(x_2)\}}
    \cdot
    d_{\mathcal{R}}\!\left(\mathcal{R}(x_1), \mathcal{R}(x_2)\right)
  \right],
\]
where the expectation is taken over pairs $(x_1, x_2)$ drawn from an
appropriate distribution over $\mathcal{X} \times \mathcal{X}$.
\end{definition}

\begin{proposition}[Zero Distortion Characterization]
Assume $d_{\mathcal{R}}(A, B) = 0$ if and only if $A = B$. Then $D_A(\pi) = 0$
if and only if $\pi$ is admissible almost surely with respect to the sampling
measure.
\end{proposition}

\begin{proof}
If $\pi$ is admissible, then whenever $\pi(x_1) = \pi(x_2)$ we have
$\mathcal{R}(x_1) = \mathcal{R}(x_2)$, so the distance term vanishes on every
pair identified by $\pi$, giving $D_A(\pi) = 0$. Conversely, if $D_A(\pi) = 0$,
the integrand is nonnegative and must vanish almost surely. By definiteness of
$d_{\mathcal{R}}$, this implies $\mathcal{R}(x_1) = \mathcal{R}(x_2)$ almost
surely on the equivalence classes of $\pi$, so $\pi$ is admissible almost
surely.
\end{proof}

For continuous representations the indicator $\mathbf{1}_{\{\pi(x_1) = \pi(x_2)\}}$
may be replaced by a smoothing kernel. Let $K_\epsilon$ be a kernel measuring
representational proximity, for example
\[
  K_\epsilon(m_1, m_2)
  =
  \exp\!\left(-\frac{d_{\mathcal{M}}(m_1, m_2)^2}{2\epsilon^2}\right).
\]
The continuous admissibility distortion is then
\[
  D_A^\epsilon(\pi)
  =
  \mathbb{E}_{x_1,\,x_2}
  \!\left[
    K_\epsilon\!\left(\pi(x_1), \pi(x_2)\right)
    \cdot
    d_{\mathcal{R}}\!\left(\mathcal{R}(x_1), \mathcal{R}(x_2)\right)
  \right],
\]
which measures whether nearby latent states possess nearby future horizons.

\begin{definition}[Lipschitz-Admissible Projection]
A projection $\pi : \mathcal{X} \to \mathcal{M}$ is $L$-Lipschitz-admissible if
\[
  d_{\mathcal{R}}\!\left(\mathcal{R}(x_1), \mathcal{R}(x_2)\right)
  \leq
  L \cdot d_{\mathcal{M}}\!\left(\pi(x_1), \pi(x_2)\right)
\]
for all $x_1, x_2 \in \mathcal{X}$.
\end{definition}

This relaxed condition requires that representational closeness must control
reachability closeness. A representation may be approximate, but it may not
place states close together in $\mathcal{M}$ while their reachable futures are
far apart.

%%% APPENDIX C
\section{Predictive Sufficiency Does Not Imply Admissibility}

Let $P(x)$ denote the predictive equivalence class of $x$: two states satisfy
$P(x_1) = P(x_2)$ when they induce the same predictions under the model and
observation process under consideration.

\begin{proposition}[Predictive--Admissibility Gap]
There exist systems for which $P(x_1) = P(x_2)$ but
$\mathcal{R}(x_1) \neq \mathcal{R}(x_2)$. Therefore predictive equivalence
does not imply admissibility equivalence.
\end{proposition}

\begin{proof}
Let $\mathcal{X} = \{a, b, c\}$ and let the observation map $o : \mathcal{X}
\to \mathcal{Y}$ satisfy $o(a) = o(b)$. Suppose the predictive model is
trained over a horizon $T$ during which the observable trajectories from $a$
and $b$ coincide, so that $P(a) = P(b)$. Define the reachability relation so
that
\[
  \mathcal{R}(a) = \{a, c\}, \qquad \mathcal{R}(b) = \{b\}.
\]
Then $c$ is reachable from $a$ but not from $b$, giving
$\mathcal{R}(a) \neq \mathcal{R}(b)$, while $P(a) = P(b)$.
\end{proof}

This construction captures the formal skeleton of the bridge-crack example.
Over an initial observation horizon, cracked and uncracked systems may behave
identically; their predictive signatures coincide. Yet their reachable futures
differ because one state admits a failure trajectory unavailable to the other.

\begin{corollary}
No predictive objective alone can guarantee admissibility unless predictive
equivalence is explicitly constrained to refine reachability equivalence.
\end{corollary}

\begin{proof}
If the predictive objective identifies $x_1$ and $x_2$ whenever
$P(x_1) = P(x_2)$, the proposition supplies cases where the induced projection
collapses states with distinct reachable futures. Such a projection is
inadmissible. An additional constraint is therefore necessary.
\end{proof}

%%% APPENDIX D
\section{Compression Sufficiency Does Not Imply Admissibility}

Let $C(x)$ denote the equivalence class induced by a compression scheme:
two states are compression-equivalent when the compressor assigns them the same
code.

\begin{proposition}[Compression--Admissibility Gap]
There exist states $x_1, x_2 \in \mathcal{X}$ such that $C(x_1) = C(x_2)$
but $\mathcal{R}(x_1) \neq \mathcal{R}(x_2)$. Therefore compression
equivalence does not imply admissibility equivalence.
\end{proposition}

\begin{proof}
Let $\mathcal{X} = \{a, b, c\}$ and suppose that $a$ and $b$ occur with
identical empirical statistics over the training window. A minimum-description
compressor assigns them the same code, since separating them does not reduce
description length over the observed sequence, giving $C(a) = C(b)$. Define
reachability so that
\[
  \mathcal{R}(a) = \{a, c\}, \qquad \mathcal{R}(b) = \{b\}.
\]
Then $a$ and $b$ possess different reachable future sets despite being
compression-equivalent. Hence compression equivalence fails to imply
admissibility equivalence.
\end{proof}

\begin{corollary}[Compression as Structural Risk]
If a compression objective rewards the collapse of empirically rare
distinctions, and if some empirically rare distinctions carry large reachability
consequences, then compression pressure can increase admissibility distortion.
\end{corollary}

\begin{proof}
A compression objective reduces code length by collapsing distinctions that do
not improve description length. If a rare distinction contributes little to
description length but separates reachable future sets, collapsing it increases
$d_{\mathcal{R}}(\mathcal{R}(x_1), \mathcal{R}(x_2))$ inside $D_A(\pi)$.
Thus compression can reduce description length while increasing admissibility
distortion.
\end{proof}

%%% APPENDIX E
\section{Optimization Inherits Projection Failure}

Let an optimizer act not on $\mathcal{X}$ directly but on a representation
$\mathcal{M}$. Let $J : \mathcal{M} \to \mathbb{R}$ be an objective function
in representation space. The implemented objective on $\mathcal{X}$ is then
$J_\pi(x) = J(\pi(x))$.

\begin{proposition}[Objective Invariance on Fibres]
For any projection $\pi : \mathcal{X} \to \mathcal{M}$, the induced objective
$J_\pi$ is constant on the fibres of $\pi$: if $\pi(x_1) = \pi(x_2)$ then
$J_\pi(x_1) = J_\pi(x_2)$.
\end{proposition}

\begin{proof}
If $\pi(x_1) = \pi(x_2)$, then
$J_\pi(x_1) = J(\pi(x_1)) = J(\pi(x_2)) = J_\pi(x_2)$.
\end{proof}

\begin{theorem}[Optimization Cannot Recover Collapsed Distinctions]
Let $\pi : \mathcal{X} \to \mathcal{M}$ be inadmissible. Then there exist
$x_1, x_2 \in \mathcal{X}$ with $\pi(x_1) = \pi(x_2)$ but
$\mathcal{R}(x_1) \neq \mathcal{R}(x_2)$. For any objective
$J : \mathcal{M} \to \mathbb{R}$, optimization of $J$ over $\mathcal{M}$
cannot distinguish $x_1$ from $x_2$.
\end{theorem}

\begin{proof}
Since $\pi$ is inadmissible, such $x_1, x_2$ exist by definition. By the
preceding proposition, $J_\pi(x_1) = J_\pi(x_2)$ for any objective $J$ on
$\mathcal{M}$. Therefore every optimizer operating through $J$ on $\mathcal{M}$
assigns the same objective value to both states, even though their reachable
future sets differ. No downstream optimization over $\mathcal{M}$ can recover
the distinction.
\end{proof}

This theorem formalizes the claim that the planner solves the wrong problem
perfectly. The failure lies not in the optimizer but in the geometry upon which
it operates.

%%% APPENDIX F
\section{Guardrail Incompleteness}

Let $G : \mathcal{M} \to \{0, 1\}$ be a safety classifier or guardrail, where
$G(m) = 0$ denotes permitted and $G(m) = 1$ denotes forbidden. Let
$\mathcal{A} : \mathcal{X} \to \{0, 1\}$ denote the true safety status of
states in the original system.

\begin{theorem}[Latent Guardrail Incompleteness]
If there exist $x_s, x_d \in \mathcal{X}$ such that $\pi(x_s) = \pi(x_d)$
but $\mathcal{A}(x_s) \neq \mathcal{A}(x_d)$, then no guardrail
$G : \mathcal{M} \to \{0, 1\}$ can correctly classify both $x_s$ and $x_d$.
\end{theorem}

\begin{proof}
Let $m = \pi(x_s) = \pi(x_d)$. Since $G$ is a function on $\mathcal{M}$, it
assigns a single value to $m$, so $G(\pi(x_s)) = G(m) = G(\pi(x_d))$. But
$\mathcal{A}(x_s) \neq \mathcal{A}(x_d)$, so at least one of the two states
must be misclassified by $G$.
\end{proof}

\begin{corollary}
A latent-space guardrail is complete only if the projection $\pi$ is admissible
with respect to the relevant safety partition.
\end{corollary}

\begin{proof}
If $\pi$ collapses any safe state and any dangerous state into the same latent
point, the theorem shows that no guardrail on $\mathcal{M}$ can classify both
correctly. Completeness therefore requires that safety-relevant distinctions be
preserved by $\pi$.
\end{proof}

%%% APPENDIX G
\section{Proxy Metrics Are Not Monotone in Admissibility}

Let $M(\pi)$ be a proxy metric associated with projection $\pi$---for instance,
predictive accuracy, compression rate, likelihood, expected reward, or
statistical fit---and let $D_A(\pi)$ be admissibility distortion. A common
implicit assumption is that improving $M$ improves admissibility. This would
require
\[
  M(\pi_1) \geq M(\pi_2)
  \quad\Longrightarrow\quad
  D_A(\pi_1) \leq D_A(\pi_2).
\]
The following proposition establishes that no such implication holds in general.

\begin{proposition}[Non-Monotonicity of Proxy Metrics]
There exist projections $\pi_1, \pi_2$ such that $M(\pi_1) > M(\pi_2)$ but
$D_A(\pi_1) > D_A(\pi_2)$. Thus improvement in a proxy metric need not reduce
admissibility distortion.
\end{proposition}

\begin{proof}
Let $\mathcal{X} = \{a, b, c\}$, and suppose $a$ and $b$ are nearly
indistinguishable according to $M$ while $c$ is easily distinguished. Let
$\pi_1$ collapse $a$ and $b$ while distinguishing $c$, and let $\pi_2$
distinguish all three states. Because $a$ and $b$ are nearly identical under
$M$, collapsing them improves the proxy metric by reducing complexity or
predictive variance, giving $M(\pi_1) > M(\pi_2)$. Now define reachability so
that $\mathcal{R}(a) \neq \mathcal{R}(b)$. Then $\pi_1$ collapses states with
different reachable futures and therefore has $D_A(\pi_1) > 0$, while $\pi_2$,
distinguishing all states, has $D_A(\pi_2) = 0$. Thus $M(\pi_1) > M(\pi_2)$
while $D_A(\pi_1) > D_A(\pi_2)$.
\end{proof}

This result captures the central formal pattern of the essay. The quantities
optimized by prediction, compression, statistics, and control are not generally
monotone in preservation of future accessibility, and there exists systematic
pressure---not merely occasional failure---for proxy improvement to trade
against admissibility near reachability-critical boundaries.

%%% APPENDIX H
\section{The Reversed Hierarchy}

The essay argues for a reversal of the usual conceptual order. The ordinary
hierarchy is
\[
  \text{Information}
  \;\rightarrow\;
  \text{Representation}
  \;\rightarrow\;
  \text{Prediction}
  \;\rightarrow\;
  \text{Planning}
  \;\rightarrow\;
  \text{Control},
\]
and the reachability hierarchy is
\[
  \text{Reachability}
  \;\rightarrow\;
  \text{Admissibility}
  \;\rightarrow\;
  \text{Representation}
  \;\rightarrow\;
  \text{Prediction}
  \;\rightarrow\;
  \text{Planning}
  \;\rightarrow\;
  \text{Control}.
\]

\begin{theorem}[Dependency of Downstream Operations]
Let a system perform prediction, planning, optimization, or control through a
projection $\pi : \mathcal{X} \to \mathcal{M}$. If $\pi$ is inadmissible,
then there exist reachability distinctions in $\mathcal{X}$ unavailable to
every downstream operation acting only on $\mathcal{M}$.
\end{theorem}

\begin{proof}
If $\pi$ is inadmissible, there exist $x_1, x_2 \in \mathcal{X}$ with
$\pi(x_1) = \pi(x_2)$ but $\mathcal{R}(x_1) \neq \mathcal{R}(x_2)$. Any
downstream operation $F$ acting only on $\mathcal{M}$ receives the same input
for both states: $F(\pi(x_1)) = F(\pi(x_2))$. Thus $F$ cannot distinguish
$x_1$ from $x_2$ despite their reachability difference. This holds for
prediction, planning, optimization, control, classification, and safety
evaluation whenever these operate only through $\mathcal{M}$.
\end{proof}

\begin{corollary}
Reachability preservation is logically prior to reliable prediction, planning,
optimization, and control under abstraction.
\end{corollary}

\begin{proof}
Each downstream operation depends upon distinctions available in the
representation. If reachability-relevant distinctions are not preserved by
the projection, no downstream operation can use them. Reliable downstream
operation therefore requires admissibility as a necessary prior condition.
\end{proof}

\bigskip

\begin{thebibliography}{99}

\bibitem{ashby1956}
Ashby, W.~Ross.
\textit{An Introduction to Cybernetics}.
Chapman \& Hall, 1956.

\bibitem{ashby1958}
Ashby, W.~Ross.
``Requisite Variety and Its Implications for the Control of Complex Systems.''
\textit{Cybernetica}, 1(2), 1958.

\bibitem{wiener1948}
Wiener, Norbert.
\textit{Cybernetics: Or Control and Communication in the Animal and the
Machine}.
MIT Press, 1948.

\bibitem{rosen1985}
Rosen, Robert.
\textit{Anticipatory Systems}.
Pergamon Press, 1985.

\bibitem{rosen1991}
Rosen, Robert.
\textit{Life Itself}.
Columbia University Press, 1991.

\bibitem{gibson1979}
Gibson, James~J.
\textit{The Ecological Approach to Visual Perception}.
Houghton Mifflin, 1979.

\bibitem{varela1991}
Varela, Francisco~J., Thompson, Evan, and Rosch, Eleanor.
\textit{The Embodied Mind}.
MIT Press, 1991.

\bibitem{shannon1948}
Shannon, Claude~E.
``A Mathematical Theory of Communication.''
\textit{Bell System Technical Journal}, 27(3):379--423, 1948.

\bibitem{kolmogorov1965}
Kolmogorov, Andrey~N.
``Three Approaches to the Quantitative Definition of Information.''
\textit{Problems of Information Transmission}, 1(1):1--7, 1965.

\bibitem{solomonoff1964}
Solomonoff, Ray~J.
``A Formal Theory of Inductive Inference.''
\textit{Information and Control}, 7(1):1--22, 1964.

\bibitem{rissanen1978}
Rissanen, Jorma.
``Modeling by Shortest Data Description.''
\textit{Automatica}, 14(5):465--471, 1978.

\bibitem{grunwald2007}
Gr\"{u}nwald, Peter~D.
\textit{The Minimum Description Length Principle}.
MIT Press, 2007.

\bibitem{schmidhuber1991}
Schmidhuber, J\"{u}rgen.
``Curious Model-Building Control Systems.''
\textit{Proceedings of the International Joint Conference on Neural Networks},
1991.

\bibitem{schmidhuber2009}
Schmidhuber, J\"{u}rgen.
``Simple Algorithmic Principles of Discovery, Subjective Beauty, Selective
Attention, Curiosity and Creativity.''
\textit{IEEE Transactions on Autonomous Mental Development}, 1(1):3--22, 2009.

\bibitem{brooks1991}
Brooks, Rodney~A.
``Intelligence Without Representation.''
\textit{Artificial Intelligence}, 47(1--3):139--159, 1991.

\bibitem{simon1969}
Simon, Herbert~A.
\textit{The Sciences of the Artificial}.
MIT Press, 1969.

\bibitem{holland1992}
Holland, John~H.
\textit{Adaptation in Natural and Artificial Systems}.
MIT Press, 1992.

\bibitem{taleb2007}
Taleb, Nassim~Nicholas.
\textit{The Black Swan}.
Random House, 2007.

\bibitem{taleb2012}
Taleb, Nassim~Nicholas.
\textit{Antifragile}.
Random House, 2012.

\bibitem{pearl2009}
Pearl, Judea.
\textit{Causality}.
Cambridge University Press, 2009.

\bibitem{friston2010}
Friston, Karl.
``The Free-Energy Principle: A Unified Brain Theory?''
\textit{Nature Reviews Neuroscience}, 11(2):127--138, 2010.

\bibitem{goodfellow2016}
Goodfellow, Ian, Bengio, Yoshua, and Courville, Aaron.
\textit{Deep Learning}.
MIT Press, 2016.

\bibitem{lecun2022}
LeCun, Yann.
``A Path Towards Autonomous Machine Intelligence.''
Open Review Position Paper, 2022.

\bibitem{bardes2022}
Bardes, Adrien, Ponce, Jean, and LeCun, Yann.
``VICReg: Variance-Invariance-Covariance Regularization for Self-Supervised
Learning.''
\textit{International Conference on Learning Representations}, 2022.

\bibitem{assran2023}
Assran, Mahmoud et al.
``Self-Supervised Learning from Images with a Joint-Embedding Predictive
Architecture.''
\textit{CVPR}, 2023.

\bibitem{park1981}
Park, David.
``Concurrency and Automata on Infinite Sequences.''
In \textit{Theoretical Computer Science}, Lecture Notes in Computer Science
vol.~104. Springer, 1981.

\bibitem{milner1989}
Milner, Robin.
\textit{Communication and Concurrency}.
Prentice Hall, 1989.

\bibitem{desharnais2004}
Desharnais, Jos\'{e}e, Gupta, Vineet, Jagadeesan, Radha, and Panangaden,
Prakash.
``Metrics for Labelled Markov Processes.''
\textit{Theoretical Computer Science}, 318(3):323--354, 2004.

\bibitem{dealfaro2004}
de Alfaro, Luca, Henzinger, Thomas~A., and Majumdar, Rupak.
``Discounting the Future in Systems Theory.''
In \textit{Automata, Languages and Programming}, ICALP 2004.

\bibitem{larsen1991}
Larsen, Kim~G. and Skou, Arne.
``Bisimulation Through Probabilistic Testing.''
\textit{Information and Computation}, 94(1):1--28, 1991.

\bibitem{dean1997}
Dean, Thomas and Givan, Robert.
``Model Minimization in Markov Decision Processes.''
\textit{Proceedings of AAAI}, 1997.

\bibitem{ladyman2007}
Ladyman, James and Ross, Don.
\textit{Every Thing Must Go: Metaphysics Naturalized}.
Oxford University Press, 2007.

\end{thebibliography}

\section{Relationship to Bisimulation and Behavioral Equivalence}

The admissibility condition---$\pi(x_1) = \pi(x_2) \Rightarrow
\mathcal{R}(x_1) = \mathcal{R}(x_2)$---will be recognized by researchers in
concurrency theory and model checking as closely related to bisimulation
equivalence. This relationship deserves explicit acknowledgment and
clarification.

Bisimulation, introduced by Park~\cite{park1981} and developed extensively
in the process algebra tradition~\cite{milner1989}, defines a relation $\sim$
between states of a transition system such that $x_1 \sim x_2$ if and only if
for every transition available from $x_1$ there is a matching transition from
$x_2$ and vice versa, with the matching states also standing in $\sim$.
Bisimilar states are behaviorally indistinguishable in the sense that no
observational experiment can separate them. Bisimulation metrics, developed
by Desharnais et al.~\cite{desharnais2004} and further by de Alfaro,
Henzinger, and Majumdar~\cite{dealfaro2004}, extend the notion to quantitative
settings by measuring the degree to which two states can be separated by any
behavioral test. The related work of Larsen and Skou~\cite{larsen1991} on
probabilistic bisimulation and Dean and Givan~\cite{dean1997} on state
abstraction in Markov decision processes are particularly relevant to the
machine learning and control applications discussed in the body of the essay.

The admissibility condition is related to these notions but differs in its
scope and motivation. Bisimulation equivalence is defined for a fixed,
fully-specified transition system and characterizes when states cannot be
distinguished by any sequence of observations within that system. Admissibility
is defined for an arbitrary projection $\pi$ out of a state space and
characterizes when the projection respects the accessibility structure relative
to a specified class of dynamics. The two conditions coincide when
$\mathcal{R}(x)$ is interpreted as the bisimulation equivalence class of $x$
and when the projection is required not to merge bisimilarly inequivalent
states.

The present framework differs from the bisimulation tradition in several
respects that are worth making explicit. First, bisimulation is typically a
property of pairs of states within a single system, while admissibility is a
property of projections out of a system into a representation space. Second,
bisimulation metrics are typically developed for probabilistic transition
systems or labeled Markov chains, while the admissibility framework applies
to any system in which reachable future sets can be defined, including
deterministic dynamical systems, set-valued dynamics, controlled systems, and
systems governed by partial or qualitative transition relations. Third, and
most importantly for the purposes of this essay, the motivation differs.
Bisimulation is primarily a tool for process equivalence, model checking, and
the verification of concurrent programs. Admissibility is proposed here as a
criterion for the evaluation of scientific abstractions, representations, and
ontologies across the full range of domains in which abstraction occurs---a
purpose for which the philosophical breadth of the structural realist tradition
\cite{ladyman2007} also provides relevant context.

Admissibility may therefore be understood as a reachability-oriented
generalization of bisimulation equivalence, extended from its home domain of
process algebra and applied to the problem of representational adequacy
across science, engineering, and machine learning. The bisimulation literature
provides rigorous foundations and a rich body of results that the present
framework can draw upon directly. The debt to that tradition is genuine and
acknowledged.

\end{document}



